\documentclass[10pt,a4paper]{amsart}
\usepackage{amsmath,amssymb,amsthm,graphicx}
% Retained minimal preamble for the Stacks Project chapter more-algebra.tex (More on Algebra).
% The chapter file pinned-original-source.tex carries only \input{preamble} on its first line, so no
% preamble text lives inside the chapter itself.  The Stacks Project uses one single shared preamble for
% every chapter; the snapshot stages it identically as sources/stacks-chow/preamble.tex and
% sources/stacks-etale-cohomology/preamble.tex.  The block below is copied mechanically, line for line,
% from sources/stacks-chow/preamble.tex lines 14-18 (the Xy-pic setup) and 154-217 (the native theorem
% declarations together with the mathematical macros).  Nothing was added, removed or redefined; the
% wrapper class and the amsmath/amssymb/amsthm/graphicx base packages are supplied by the assembler.
% Chapter: more-algebra.tex -- More on Algebra.
% For commutative diagrams we use Xy-pic
\usepackage[all]{xy}

% We use 2cell for 2-commutative diagrams.
\xyoption{2cell}

% Theorem environments.
%
\theoremstyle{plain}
\newtheorem{theorem}[subsection]{Theorem}
\newtheorem{proposition}[subsection]{Proposition}
\newtheorem{lemma}[subsection]{Lemma}

\theoremstyle{definition}
\newtheorem{definition}[subsection]{Definition}
\newtheorem{example}[subsection]{Example}
\newtheorem{exercise}[subsection]{Exercise}
\newtheorem{situation}[subsection]{Situation}

\theoremstyle{remark}
\newtheorem{remark}[subsection]{Remark}
\newtheorem{remarks}[subsection]{Remarks}

\numberwithin{equation}{subsection}

% Macros
%
\def\lim{\mathop{\mathrm{lim}}\nolimits}
\def\colim{\mathop{\mathrm{colim}}\nolimits}
\def\Spec{\mathop{\mathrm{Spec}}}
\def\Hom{\mathop{\mathrm{Hom}}\nolimits}
\def\Ext{\mathop{\mathrm{Ext}}\nolimits}
\def\SheafHom{\mathop{\mathcal{H}\!\mathit{om}}\nolimits}
\def\SheafExt{\mathop{\mathcal{E}\!\mathit{xt}}\nolimits}
\def\Sch{\mathit{Sch}}
\def\Mor{\mathop{\mathrm{Mor}}\nolimits}
\def\Ob{\mathop{\mathrm{Ob}}\nolimits}
\def\Sh{\mathop{\mathit{Sh}}\nolimits}
\def\NL{\mathop{N\!L}\nolimits}
\def\CH{\mathop{\mathrm{CH}}\nolimits}
\def\proetale{{pro\text{-}\acute{e}tale}}
\def\etale{{\acute{e}tale}}
\def\QCoh{\mathit{QCoh}}
\def\Ker{\mathop{\mathrm{Ker}}}
\def\Im{\mathop{\mathrm{Im}}}
\def\Coker{\mathop{\mathrm{Coker}}}
\def\Coim{\mathop{\mathrm{Coim}}}

% Boxtimes
%
\DeclareMathSymbol{\boxtimes}{\mathbin}{AMSa}{"02}

%
% Macros for moduli stacks/spaces
%
\def\QCohstack{\mathcal{QC}\!\mathit{oh}}
\def\Cohstack{\mathcal{C}\!\mathit{oh}}
\def\Spacesstack{\mathcal{S}\!\mathit{paces}}
\def\Quotfunctor{\mathrm{Quot}}
\def\Hilbfunctor{\mathrm{Hilb}}
\def\Curvesstack{\mathcal{C}\!\mathit{urves}}
\def\Polarizedstack{\mathcal{P}\!\mathit{olarized}}
\def\Complexesstack{\mathcal{C}\!\mathit{omplexes}}
% \Pic is the operator that assigns to X its picard group, usage \Pic(X)
% \Picardstack_{X/B} denotes the Picard stack of X over B
% \Picardfunctor_{X/B} denotes the Picard functor of X over B
\def\Pic{\mathop{\mathrm{Pic}}\nolimits}
\def\Picardstack{\mathcal{P}\!\mathit{ic}}
\def\Picardfunctor{\mathrm{Pic}}
\def\Deformationcategory{\mathcal{D}\!\mathit{ef}}

% --- appended 2026-10-12 (Chief): the project own definitions of the reference/slogan/history
% display environments, copied mechanically from the shared Stacks preamble lines 8-12
% (\usepackage{verbatim} plus the three \newenvironment lines, which the chapter preamble keeps
% verbatim from the project preamble). Nothing added, removed or redefined: these environments
% carry no mathematical content and the project itself typesets them as comments. Without these
% three lines a counted statement whose own text uses them fails with Environment slogan undefined.
\usepackage{verbatim}
\newenvironment{reference}{\comment}{\endcomment}
\newenvironment{slogan}{\comment}{\endcomment}
\newenvironment{history}{\comment}{\endcomment}
\usepackage{cleveref}
\crefname{originallemma}{original lemma}{original lemmas}
\Crefname{originallemma}{Original lemma}{Original lemmas}
\makeatletter
\newlabel{lemma-base-change-weakly-etale}{{092H}{source}{}{sourcepointer-1}{}}
\newlabel{lemma-base-change-weakly-etale@cref}{{[originallemma][0][]092H}{[1][0][]source}{}{}{}}
\makeatother
\crefname{originallemma}{original lemma}{original lemmas}
\Crefname{originallemma}{Original lemma}{Original lemmas}
\makeatletter
\newlabel{lemma-max-weakly-etale-subalgebra}{{0CKS}{source}{}{sourcepointer-2}{}}
\newlabel{lemma-max-weakly-etale-subalgebra@cref}{{[originallemma][0][]0CKS}{[1][0][]source}{}{}{}}
\makeatother
\crefname{originallemma}{original lemma}{original lemmas}
\Crefname{originallemma}{Original lemma}{Original lemmas}
\makeatletter
\newlabel{lemma-properties-of-max-weakly-etale-subalgebra}{{0CKT}{source}{}{sourcepointer-3}{}}
\newlabel{lemma-properties-of-max-weakly-etale-subalgebra@cref}{{[originallemma][0][]0CKT}{[1][0][]source}{}{}{}}
\makeatother
\crefname{originallemma}{original lemma}{original lemmas}
\Crefname{originallemma}{Original lemma}{Original lemmas}
\makeatletter
\newlabel{lemma-class-weakly-etale-over-field}{{0CKR}{source}{}{sourcepointer-4}{}}
\newlabel{lemma-class-weakly-etale-over-field@cref}{{[originallemma][0][]0CKR}{[1][0][]source}{}{}{}}
\makeatother
\crefname{originallemma}{original lemma}{original lemmas}
\Crefname{originallemma}{Original lemma}{Original lemmas}
\makeatletter
\newlabel{lemma-weakly-etale-permanence}{{092L}{source}{}{sourcepointer-5}{}}
\newlabel{lemma-weakly-etale-permanence@cref}{{[originallemma][0][]092L}{[1][0][]source}{}{}{}}
\makeatother
\crefname{originallemma}{original lemma}{original lemmas}
\Crefname{originallemma}{Original lemma}{Original lemmas}
\makeatletter
\newlabel{algebra-lemma-Noetherian-irreducible-components}{{00FR}{source}{}{sourcepointer-6}{}}
\newlabel{algebra-lemma-Noetherian-irreducible-components@cref}{{[originallemma][0][]00FR}{[1][0][]source}{}{}{}}
\makeatother
\crefname{originallemma}{original lemma}{original lemmas}
\Crefname{originallemma}{Original lemma}{Original lemmas}
\makeatletter
\newlabel{algebra-lemma-Noetherian-permanence}{{00FN}{source}{}{sourcepointer-7}{}}
\newlabel{algebra-lemma-Noetherian-permanence@cref}{{[originallemma][0][]00FN}{[1][0][]source}{}{}{}}
\makeatother
\crefname{originallemma}{original lemma}{original lemmas}
\Crefname{originallemma}{Original lemma}{Original lemmas}
\makeatletter
\newlabel{algebra-lemma-make-geometrically-irreducible}{{037Q}{source}{}{sourcepointer-8}{}}
\newlabel{algebra-lemma-make-geometrically-irreducible@cref}{{[originallemma][0][]037Q}{[1][0][]source}{}{}{}}
\makeatother
\crefname{originallemma}{original lemma}{original lemmas}
\Crefname{originallemma}{Original lemma}{Original lemmas}
\makeatletter
\newlabel{algebra-lemma-minimal-prime-reduced-ring}{{00EU}{source}{}{sourcepointer-9}{}}
\newlabel{algebra-lemma-minimal-prime-reduced-ring@cref}{{[originallemma][0][]00EU}{[1][0][]source}{}{}{}}
\makeatother
\crefname{originallemma}{original lemma}{original lemmas}
\Crefname{originallemma}{Original lemma}{Original lemmas}
\makeatletter
\newlabel{algebra-lemma-total-ring-fractions-no-embedded-points}{{02LX}{source}{}{sourcepointer-10}{}}
\newlabel{algebra-lemma-total-ring-fractions-no-embedded-points@cref}{{[originallemma][0][]02LX}{[1][0][]source}{}{}{}}
\makeatother
\crefname{originallemma}{original lemma}{original lemmas}
\Crefname{originallemma}{Original lemma}{Original lemmas}
\makeatletter
\newlabel{fields-lemma-finite-separable-tensor-alg-closed}{{0CKN}{source}{}{sourcepointer-11}{}}
\newlabel{fields-lemma-finite-separable-tensor-alg-closed@cref}{{[originallemma][0][]0CKN}{[1][0][]source}{}{}{}}
\makeatother
\title{The maximal weakly etale subalgebra commutes with an extension of the base field}
\author{The Stacks Project Authors}
\date{stacks-project:more-algebra.tex:local-snapshot-20261012; GNU FDL 1.2 or later}
\begin{document}
\maketitle
\setcounter{subsection}{1}
\noindent\textit{Editorial note.} All delivered mathematical text is the author's own, copied from the pinned Stacks Project chapter file more-algebra.tex; no statement or proof step was rewritten, reordered or removed. The counted span is the author environment and its own proof environment. References inside the span are handed to the assembler as navigation stubs that carry the canonical tag of the cited tag only; whether each of them is load-bearing is recorded on the stub itself.\par\medskip
\noindent\textit{Allowed standard prerequisites (editorial).}\begin{itemize}
\item The maximal weakly etale subalgebra of an algebra over a field is weakly etale, is preserved by base change of the field, and is a field when the algebra is reduced and geometrically irreducible. invoked at proof lines 30988-30994 and again at proof lines 31064-31072. Identifies the two sides of the comparison map and supplies the field structure used in the final computation.
\item Base change of a weakly etale algebra along a field extension is weakly etale, which makes the comparison map injective. invoked at proof lines 30991-30993. Produces the canonical injective map from the base changed subalgebra into the maximal subalgebra of the base changed algebra.
\item The maximal weakly etale subalgebra is compatible with filtered colimits, so a statement about finite type algebras implies the general case. invoked at proof lines 31003-31006. Reduces the proof to the case of a finite type algebra over the base field.
\item Weakly etale algebras over a field are reduced, which lets the nilpotent kernel of the reduction be discarded. invoked at proof lines 31010-31014. Shows that the comparison map stays injective when the algebra is replaced by its reduced quotient.
\item A finite separable extension of fields becomes a product of copies of the larger field after a base change to a separably closed or algebraically closed field. invoked at proof lines 31083-31088. Decomposes the tensor product of the algebra over the finite separable field with the algebraically closed field into a product of factors.
\end{itemize}\medskip
\ifdefined\hypertarget\hypertarget{sourcepointer-1}{}\fi
\noindent\textit{Source reference.} Original lemma 092H: sources/stacks-more-algebra/pinned-original-source.tex, source lines 30303--30303. Used at proof lines 30991-30993 to see that the base changed maximal subalgebra is weakly etale over the larger field, which produces the comparison map; everything built on that map is inside the counted span.\par\medskip
\ifdefined\hypertarget\hypertarget{sourcepointer-2}{}\fi
\noindent\textit{Source reference.} Original lemma 0CKS: sources/stacks-more-algebra/pinned-original-source.tex, source lines 30892--30892. The statement names the object through this lemma; it is a pointer fixing notation, and every property of that object used later is invoked explicitly and applied inside the counted span.\par\medskip
\ifdefined\hypertarget\hypertarget{sourcepointer-3}{}\fi
\noindent\textit{Source reference.} Original lemma 0CKT: sources/stacks-more-algebra/pinned-original-source.tex, source lines 30913--30913. Invoked for compatibility with filtered colimits at proof lines 31003-31006, for the behaviour under reduction at proof lines 31013-31016 and for the field and geometric irreducibility properties at proof lines 31064-31072; each application and the computation that follows are inside the counted span.\par\medskip
\ifdefined\hypertarget\hypertarget{sourcepointer-4}{}\fi
\noindent\textit{Source reference.} Original lemma 0CKR: sources/stacks-more-algebra/pinned-original-source.tex, source lines 30846--30846. Used at proof lines 31010-31012 to conclude that the maximal subalgebra over the larger field is reduced, which is what makes the reduction step work; the reduction itself is inside the counted span.\par\medskip
\ifdefined\hypertarget\hypertarget{sourcepointer-5}{}\fi
\noindent\textit{Source reference.} Original lemma 092L: sources/stacks-more-algebra/pinned-original-source.tex, source lines 30405--30405. The final sentence at proof lines 31098-31101 uses it to name the map whose weak etaleness gives the equality; the decomposition producing that map is inside the counted span.\par\medskip
\ifdefined\hypertarget\hypertarget{sourcepointer-6}{}\fi
\noindent\textit{Source reference.} Original lemma 00FR: sources/stacks-chow/tags.txt, source lines 1--1. A navigation pointer for the decompositions of spectra into irreducible components used while reducing to the reduced case; the use made of the decomposition is inside the counted span.\par\medskip
\ifdefined\hypertarget\hypertarget{sourcepointer-7}{}\fi
\noindent\textit{Source reference.} Original lemma 00FN: sources/stacks-chow/tags.txt, source lines 1--1. A navigation pointer for the permanence statements the reduction to finite type algebras relies on; the reduction itself and its consequences are inside the counted span.\par\medskip
\ifdefined\hypertarget\hypertarget{sourcepointer-8}{}\fi
\noindent\textit{Source reference.} Original lemma 037Q: sources/stacks-chow/tags.txt, source lines 1--1. Used at proof lines 31071-31072 to see that the algebra is geometrically irreducible over the finite separable field of constants; the decomposition that uses this is inside the counted span.\par\medskip
\ifdefined\hypertarget\hypertarget{sourcepointer-9}{}\fi
\noindent\textit{Source reference.} Original lemma 00EU: sources/stacks-chow/tags.txt, source lines 1--1. A navigation pointer used during the reduction to the reduced case; the comparison map and the reduction step are inside the counted span.\par\medskip
\ifdefined\hypertarget\hypertarget{sourcepointer-10}{}\fi
\noindent\textit{Source reference.} Original lemma 02LX: sources/stacks-chow/tags.txt, source lines 1--1. A navigation pointer for the manipulation of the field of constants in the reduced case; the identification of that field is inside the counted span.\par\medskip
\ifdefined\hypertarget\hypertarget{sourcepointer-11}{}\fi
\noindent\textit{Source reference.} Original lemma 0CKN: sources/stacks-chow/tags.txt, source lines 1--1. Invoked at proof lines 31085-31090 for the displayed isomorphism splitting the tensor product into copies of the algebraically closed field; the subsequent factor by factor computation of the maximal subalgebras is inside the counted span.\par\medskip
\begin{lemma}
\label{lemma-change-fields-max-weakly-etale-subalgebra}
Let $L/K$ be an extension of fields. Let $A$ be a $K$-algebra.
Let $B \subset A$ be the maximal weakly \'etale $K$-subalgebra of
$A$ as in Lemma \ref{lemma-max-weakly-etale-subalgebra}.
Then $B \otimes_K L$ is the maximal weakly \'etale $L$-subalgebra
of $A \otimes_K L$.
\end{lemma}

\begin{proof}
For an algebra $A$ over $K$ we write $B_{max}(A/K)$
for the maximal weakly \'etale $K$-subalgebra of $A$.
Similarly we write $B_{max}(A'/L)$ for the maximal weakly \'etale
$L$-subalgebra of $A'$ if $A'$ is an $L$-algebra.
Since $B_{max}(A/K) \otimes_K L$ is weakly \'etale over $L$
(Lemma \ref{lemma-base-change-weakly-etale})
and since $B_{max}(A/K) \otimes_K L \subset A \otimes_K L$
we obtain a canonical injective map
$$
B_{max}(A/K) \otimes_K L \to B_{max}((A \otimes_K L)/L)
$$
The lemma states that this map is an isomorphism.

\medskip\noindent
To prove the lemma for $L$ and our $K$-algebra $A$, it suffices to
prove the lemma for any field extension $L'$ of $L$. Namely, we have
the factorization
$$
B_{max}(A/K) \otimes_K L' \to
B_{max}((A \otimes_K L)/L) \otimes_L L' \to
B_{max}((A \otimes_K L')/L')
$$
hence the composition cannot be surjective without
$B_{max}(A/K) \otimes_K L \to B_{max}((A \otimes_K L)/L)$
being surjective. Thus we may assume $L$ is algebraically closed.

\medskip\noindent
Reduction to finite type $K$-algebra. We may write $A$
is the filtered colimit of its finite type $K$-subalgebras.
Using Lemma \ref{lemma-properties-of-max-weakly-etale-subalgebra}
we see that it suffices to prove the lemma for finite type
$K$-algebras.

\medskip\noindent
Assume $A$ is a finite type $K$-algebra. Since the kernel
of $A \to A_{red}$ is nilpotent, the same is true for
$A \otimes_K L \to A_{red} \otimes_K L$. Then
$$
B_{max}((A \otimes_K L)/L) \to B_{max}((A_{red} \otimes_K L)/L)
$$
is injective because the kernel is nilpotent and the
weakly \'etale $L$-algebra $B_{max}((A \otimes_K L)/L)$ is reduced
(Lemma \ref{lemma-class-weakly-etale-over-field}).
Since $B_{max}(A/K) = B_{max}(A_{red}/K)$
by Lemma \ref{lemma-properties-of-max-weakly-etale-subalgebra}
we conclude
that it suffices to prove the lemma for $A_{red}$.

\medskip\noindent
Assume $A$ is a reduced finite type $K$-algebra.
Let $Q = Q(A)$ be the total quotient ring of $A$.
Then $A \subset Q$ and $A \otimes_K L \subset Q \otimes_A L$
and hence
$$
B_{max}(A/K) = A \cap B_{max}(Q/K)
$$
and
$$
B_{max}((A \otimes_K L)/L) =
(A \otimes_K L) \cap B_{max}((Q \otimes_K L)/L)
$$
by Lemma \ref{lemma-properties-of-max-weakly-etale-subalgebra}.
Since $-\otimes_K L$ is an exact functor, it follows that
if we prove the result for $Q$, then the result follows for $A$.
Since $Q$ is a finite product of fields (Algebra, Lemmas
\ref{algebra-lemma-total-ring-fractions-no-embedded-points},
\ref{algebra-lemma-minimal-prime-reduced-ring},
\ref{algebra-lemma-Noetherian-irreducible-components}, and
\ref{algebra-lemma-Noetherian-permanence})
and since $B_{max}$ commutes with products
(Lemma \ref{lemma-properties-of-max-weakly-etale-subalgebra})
it suffices to prove the lemma when $A$ is a field.

\medskip\noindent
Assume $A$ is a field. We reduce to $A$ being finitely generated
over $K$ by the argument in the third paragraph of the proof.
(In fact the way we reduced to the case of a field produces
a finitely generated field extension of $K$.)

\medskip\noindent
Assume $A$ is a finitely generated field extension of $K$.
Then $K' = B_{max}(A/K)$ is a field separable algebraic over $K$ by
Lemma \ref{lemma-properties-of-max-weakly-etale-subalgebra} part (6).
Hence $K'$ is a finite separable field extension
of $K$ and $A$ is geometrically irreducible over $K'$ by
Algebra, Lemma \ref{algebra-lemma-make-geometrically-irreducible}.
Since $L$ is algebraically closed and $K'/K$ finite separable
we see that
$$
K' \otimes_K L \to \prod\nolimits_{\sigma \in \Hom_K(K', L)} L,\quad
\alpha \otimes \beta \mapsto (\sigma(\alpha)\beta)_\sigma
$$
is an isomorphism
(Fields, Lemma \ref{fields-lemma-finite-separable-tensor-alg-closed}).
We conclude
$$
A \otimes_K L = A \otimes_{K'} (K' \otimes_K L) =
\prod\nolimits_{\sigma \in \Hom_K(K', L)} A \otimes_{K', \sigma} L
$$
Since $A$ is geometrically irreducible over $K'$ we see that
$A \otimes_{K', \sigma} L$ has a unique minimal prime.
Since $L$ is algebraically closed it follows that
$B_{max}((A \otimes_{K', \sigma} L)/L) = L$
because this $L$-algebra is a field algebraic over $L$ by
Lemma \ref{lemma-properties-of-max-weakly-etale-subalgebra} part (6).
It follows that the maximal weakly \'etale $K' \otimes_K L$-subalgebra
of $A \otimes_K L$ is $K' \otimes_K L$ because we can decompose
these subalgebras into products as above. Hence the inclusion
$K' \otimes_K L \subset B_{max}((A \otimes_K L)/L)$ is an
equality: the ring map $K' \otimes_K L \to B_{max}((A \otimes_K L)/L)$
is weakly \'etale by
Lemma \ref{lemma-weakly-etale-permanence}.
\end{proof}

\end{document}
